\documentclass[11pt,a4paper]{article}

\usepackage[margin=2.2cm]{geometry}
\usepackage[T1]{fontenc}
\usepackage[utf8]{inputenc}
\usepackage{lmodern}
\usepackage{microtype}
\usepackage{amsmath,amssymb}
\usepackage{enumitem}
\usepackage{xcolor}
\usepackage{listings}
\usepackage{hyperref}
\usepackage{titlesec}
\usepackage{fancyhdr}

\definecolor{codegray}{gray}{0.96}
\definecolor{linkblue}{RGB}{34,83,145}
\hypersetup{colorlinks=true,linkcolor=linkblue,urlcolor=linkblue}

\lstset{
  basicstyle=\ttfamily\small,
  backgroundcolor=\color{codegray},
  frame=single,
  rulecolor=\color{gray!30},
  breaklines=true,
  columns=fullflexible,
  showstringspaces=false,
  xleftmargin=0.3em,
  xrightmargin=0.3em
}

\titleformat{\section}{\Large\bfseries}{\thesection.}{0.5em}{}
\titleformat{\subsection}{\large\bfseries}{\thesubsection.}{0.5em}{}
\setlist[itemize]{topsep=0.3em,itemsep=0.25em}
\setlist[enumerate]{topsep=0.3em,itemsep=0.25em}

\pagestyle{fancy}
\fancyhf{}
\lhead{Generative Models -- M2 Computer Science}
\rhead{VAE Practical Lab}
\cfoot{\thepage}
\setlength{\headheight}{14pt}

\newcommand{\question}[2]{%
  \vspace{0.4em}
  \noindent\textbf{Question #1.} #2
  \vspace{0.5em}
}

\begin{document}

\begin{center}
  {\LARGE\bfseries Practical Lab: From Autoencoders to Variational Autoencoders}\par
  \vspace{0.5em}
  {\large Generative Models -- M2 Computer Science}\par
  \vspace{0.8em}
  \begin{tabular}{ll}
    \textbf{Duration:} & 2 hours, including setup and a short report \\
    \textbf{Dataset:} & MNIST \\
    \textbf{Hardware:} & one small GPU per student
  \end{tabular}
\end{center}

\vspace{0.8em}

\section*{Learning goals}

By the end of this lab, you should be able to:
\begin{enumerate}
  \item train and use a simple autoencoder (AE);
  \item visualize a 2D latent representation;
  \item explain why a good autoencoder is \textbf{not automatically a good generative model};
  \item fit a Gaussian mixture model (GMM) to the latent codes of an autoencoder and sample from it;
  \item implement the three key ingredients of a variational autoencoder (VAE):
  \begin{itemize}
    \item a Gaussian encoder,
    \item the reparameterization trick,
    \item the reconstruction + KL objective;
  \end{itemize}
  \item generate new MNIST digits and explore the learned latent space.
\end{enumerate}

The lab is deliberately small. Most of the PyTorch infrastructure is already written. Your task is to complete a few important lines and interpret the results.

\section{Setup}

Clone the repository provided by the instructor, then create the environment:

\begin{lstlisting}[language=bash]
conda env create -f environment.yml
conda activate genmodels-vae
\end{lstlisting}

Check that PyTorch, CUDA, and the dataset are available:

\begin{lstlisting}[language=bash]
python check_install.py
\end{lstlisting}

If MNIST is stored in a shared directory on the cluster, use:

\begin{lstlisting}[language=bash]
python check_install.py --data-dir /path/to/shared/MNIST
\end{lstlisting}

The expected final lines are similar to:

\begin{lstlisting}
PyTorch: OK
CUDA: OK
MNIST: OK
Everything is ready!
\end{lstlisting}

\noindent\textbf{Important.} The practical work is in \texttt{student.py}. Search for \texttt{TODO}. You should not need to modify \texttt{run\_lab.py} or \texttt{utils.py}.

You can list all TODOs with:

\begin{lstlisting}[language=bash]
grep -n "TODO" student.py
\end{lstlisting}

\section{Autoencoder: good reconstruction does not imply good generation}

An autoencoder contains an encoder $E_\phi$ and a decoder $D_\theta$:
\[
z = E_\phi(x), \qquad \hat x = D_\theta(z).
\]
We train it to reconstruct the input:
\[
\min_{\theta,\phi} \sum_i \|x_i - D_\theta(E_\phi(x_i))\|^2.
\]
In this lab the latent variable has dimension \textbf{2}, so we can visualize it directly.

\subsection*{Task 1 -- Complete the autoencoder forward pass}

Open \texttt{student.py} and complete \textbf{TODO 1}.

Then train the autoencoder:

\begin{lstlisting}[language=bash]
python run_lab.py ae
\end{lstlisting}

The script creates:
\begin{itemize}
  \item \texttt{outputs/ae\_reconstructions.png}
  \item \texttt{outputs/ae\_latent.png}
  \item \texttt{outputs/ae\_prior\_samples.png}
  \item \texttt{outputs/ae.pt}
\end{itemize}

Look at all three figures before continuing.

\question{1}{The autoencoder can reconstruct digits, but samples obtained by drawing
\[
z\sim\mathcal N(0,I_2)
\]
and decoding them are often poor. Why is this not surprising?}

\section{Learning a distribution in the latent space with a GMM}

The autoencoder has learned a set of latent codes $z_i=E_\phi(x_i)$, but it has \textbf{not} learned a probability distribution from which we can sample latent variables.

A simple idea is therefore:
\begin{enumerate}
  \item encode the MNIST training images;
  \item fit a Gaussian mixture model to their 2D latent codes;
  \item sample new latent variables from the GMM;
  \item decode them with the already-trained decoder.
\end{enumerate}

A mixture of $K$ Gaussians has density
\[
p(z)=\sum_{k=1}^K \pi_k\,\mathcal N(z;\mu_k,\Sigma_k).
\]
We use \texttt{sklearn.mixture.GaussianMixture}, which performs the EM optimization for us.

\subsection*{Task 2 -- Fit and sample the GMM}

Complete \textbf{TODO 2} and \textbf{TODO 3} in \texttt{student.py}.

Then run:

\begin{lstlisting}[language=bash]
python run_lab.py gmm
\end{lstlisting}

This creates:
\begin{itemize}
  \item \texttt{outputs/ae\_latent\_gmm.png}
  \item \texttt{outputs/ae\_gmm\_samples.png}
\end{itemize}

\question{2}{Compare \texttt{ae\_prior\_samples.png} and \texttt{ae\_gmm\_samples.png}. Why should sampling from the fitted GMM work better than sampling from $\mathcal N(0,I_2)$?}

\question{3}{After fitting the GMM, we can generate according to
\[
z\sim p_{\mathrm{GMM}}(z), \qquad x=D_\theta(z).
\]
In what sense is \textbf{GMM + decoder} now a generative model?}

\section{Variational autoencoder}

Instead of fitting a distribution to the latent codes \textbf{after} training, a VAE learns a latent representation that is compatible with a chosen prior during training.

We use the generative model from the course:
\[
z\sim\mathcal N(0,I_2),
\qquad
x\mid z\sim\mathcal N(D_\theta(z),I).
\]

The exact posterior $p_\theta(z\mid x)$ is difficult to compute, so the encoder represents an approximate Gaussian posterior
\[
q_\phi(z\mid x)
=\mathcal N\!\left(\mu_\phi(x),\operatorname{diag}(\sigma_\phi(x)^2)\right).
\]
The network will output \texttt{mu} and \texttt{logvar = log(sigma\^{}2)}.

\subsection*{Task 3 -- Gaussian encoder}

Complete \textbf{TODO 4} so that the encoder returns the mean and log-variance of the approximate posterior.

\subsection*{Task 4 -- Reparameterization trick}

To backpropagate through a random latent variable, write
\[
\varepsilon\sim\mathcal N(0,I),
\qquad
z=\mu+\sigma\odot\varepsilon.
\]
Complete \textbf{TODO 5}.

Useful PyTorch functions include:

\begin{lstlisting}[language=Python]
torch.exp(...)
torch.randn_like(...)
\end{lstlisting}

Remember that the network outputs \texttt{logvar = log(sigma\^{}2)}.

\subsection*{Task 5 -- VAE objective}

For a diagonal Gaussian posterior and a standard Gaussian prior,
\[
D_{\mathrm{KL}}\!\left(
\mathcal N(\mu,\operatorname{diag}(\sigma^2))
\,\|\,
\mathcal N(0,I)
\right)
=
\frac12\sum_j\left(\mu_j^2+\sigma_j^2-1-\log\sigma_j^2\right).
\]

Because the decoder likelihood is Gaussian, the reconstruction contribution is proportional to a squared reconstruction error.

We minimize
\[
\mathcal J
=
\underbrace{\|x-D_\theta(z)\|^2}_{\text{reconstruction}}
+\beta\,
\underbrace{D_{\mathrm{KL}}(q_\phi(z\mid x)\|\mathcal N(0,I))}_{\text{regularization}}.
\]

Complete \textbf{TODO 6} and \textbf{TODO 7} in \texttt{student.py}.

Then train the VAE:

\begin{lstlisting}[language=bash]
python run_lab.py vae
\end{lstlisting}

The script creates:
\begin{itemize}
  \item \texttt{outputs/vae\_reconstructions.png}
  \item \texttt{outputs/vae\_latent.png}
  \item \texttt{outputs/vae\_prior\_samples.png}
  \item \texttt{outputs/vae\_training\_curves.png}
  \item \texttt{outputs/vae.pt}
\end{itemize}

\question{4}{What is the role of the KL term? Compare the AE and VAE latent spaces.}

\question{5}{Why do we use the reparameterization
\[
z=\mu+\sigma\varepsilon
\]
instead of simply writing ``sample $z\sim q_\phi(z\mid x)$'' inside the network?}

\section{Explore the VAE latent space}

Now use the trained VAE as a generative model.

Run:

\begin{lstlisting}[language=bash]
python run_lab.py explore
\end{lstlisting}

This produces:
\begin{itemize}
  \item \texttt{outputs/vae\_latent\_grid.png}: decode a regular grid of points in latent space;
  \item \texttt{outputs/vae\_interpolation.png}: encode two MNIST images and interpolate between their latent means.
\end{itemize}

\question{6}{Look at the latent grid and the interpolation. Describe one qualitative property of the latent representation that is useful for generation.}

\section{Optional extension -- what does beta do?}

If you finish early, retrain the VAE with a different weight on the KL term:

\begin{lstlisting}[language=bash]
python run_lab.py vae --beta 0 --tag beta0
python run_lab.py vae --beta 5 --tag beta5
\end{lstlisting}

You can add \texttt{--tag beta0}, \texttt{--tag beta5}, etc. to keep several runs without overwriting previous figures.

Compare reconstruction quality, latent organization, and samples from the prior.

\noindent\textbf{Optional question.} What trade-off do you observe when $\beta$ is increased?

\section{Report}

Submit a \textbf{short report (maximum 2 pages)}. A LaTeX template is provided in \path{report_template.tex}.

Your report should contain:
\begin{enumerate}
  \item one AE reconstruction figure and a short comment;
  \item a comparison of AE samples from $\mathcal N(0,I)$ and from the fitted GMM;
  \item one VAE latent-space figure and one VAE generation figure;
  \item concise answers to Questions 1--6.
\end{enumerate}

We are interested in your \textbf{interpretation}, not in a description of every line of code.

\section*{Useful commands}

Fast debugging run on a small subset:

\begin{lstlisting}[language=bash]
python run_lab.py ae --quick
python run_lab.py vae --quick
\end{lstlisting}

Choose CPU explicitly:

\begin{lstlisting}[language=bash]
python run_lab.py ae --device cpu
\end{lstlisting}

Choose another data directory:

\begin{lstlisting}[language=bash]
python run_lab.py ae --data-dir /path/to/MNIST
\end{lstlisting}

Clean generated outputs:

\begin{lstlisting}[language=bash]
rm -rf outputs
\end{lstlisting}

\vfill
\begin{center}
\textit{Good luck.}
\end{center}

\end{document}
